\documentclass{article}
\usepackage{tcee-repaso}
\tcetema[Unión Bancaria y Unión de Mercados de Capitales]{3.B.44}{La Unión Bancaria: Pilares y Código normativo único. La Unión de Mercados de Capitales.}
\tcefecha{14/11/2026}
\begin{document}
\cabecerarepaso
\begin{columnas}

% ==========================================================================
\bloque{Esquema del cante}
\esquemacante{%
\apartadocante{Introducción: «empezar la casa por el tejado»}{2,5}
\apartadocante{I. Origen, fundamentos y código normativo único}{7,5}
\subcante{1. Origen: de la crisis a la decisión de 2012}{2}
\subcante{2. ¿Por qué? Círculo vicioso y trilema financiero}{3,5}
\subcante{3. El código normativo único}{2}
\apartadocante{II. Los pilares de la unión bancaria}{10,5}
\subcante{1. Mecanismo Único de Supervisión}{3}
\subcante{2. Mecanismo Único de Resolución, con el \emph{bail-in}}{4}
\subcante{3. La garantía común de depósitos, con el seguro común}{3,5}
\apartadocante{III. La Unión de Mercados de Capitales}{8,5}
\subcante{1. ¿Por qué?, con el reparto de riesgos}{3}
\subcante{2. La UMC (2015-2024)}{2,5}
\subcante{3. La Unión de Ahorros e Inversiones (2025-2026)}{3}
\apartadocante{Conclusión}{1}}
\emph{Gráficos y figuras del cante}: círculo vicioso en $(L,s)$ y trilema financiero (\hyperref[rep:I]{I}); pilares, solvencia, FUR, jerarquía de acreedores y fondo común (\hyperref[rep:II]{II}); financiación de las empresas, capitalización y reparto de riesgos (\hyperref[rep:III]{III}). Fórmulas: $s^*=\frac{s_0+\alpha L_0}{1-\alpha\beta}$; $\theta B\geq C$ frente a $B\geq C$; $r-r_f\approx\pi(1-R)$; $F_c/F_n=\sqrt{\rho+(1-\rho)/n}$; $1=\beta_K+\beta_F+\beta_C+\beta_U$.

% ==========================================================================
\bloque[2,5]{Introducción}
\begin{itemize}
\item \textbf{Enganche}: «Empezar la casa por el tejado»: en 1999, moneda y banco central comunes, pero bancos supervisados, rescatados y garantizados por cada Estado. Irlanda gastó más de un tercio de su PIB en sus bancos. ¿Puede una unión monetaria funcionar sin un sistema financiero común?
\item \textbf{Relevancia}: el BCE supervisa 109 grupos, $\approx$85\,\% de los activos, $\approx$29,5 billones; préstamos = 73\,\% de la deuda de las empresas (EE.\,UU.: 37\,\%); \autor{Draghi}: 750.000-800.000 millones más de inversión al año. España: BdE en el MUS, FROB en el MUR, Popular (2017) \vertema{4.A.19} \vertema{4.A.20}.
\item \textbf{Contextualización}: \concepto{unión bancaria}: supervisión, resolución (y garantía) europeas sobre el \concepto{código normativo único}; euro + cooperación estrecha. \concepto{UMC} (2015) $\to$ \concepto{UAI} (2025): los 27 \vertema{3.B.39} \vertema{3.B.43} \vertema{3.B.45}.
\item \textbf{Problemática}: ¿por qué en 2012? ¿Qué normas y pilares? ¿Qué falta? ¿Por qué mercados integrados?
\item \textbf{Estructura}: \hyperref[rep:I]{I} porqué y normas; \hyperref[rep:II]{II} pilares; \hyperref[rep:III]{III} mercados.
\end{itemize}

% ==========================================================================
\bloque[7,5]{I. Origen, fundamentos y código normativo único\label{rep:I}}
\begin{ideaclave}
Con bancos nacionales, los problemas de los bancos y del Estado se alimentan entre sí, y un supervisor nacional no valora el coste para los vecinos. Respuesta: supervisión y resolución europeas sobre normas comunes para los 27.
\end{ideaclave}

\sub{I.1. Origen\label{rep:I.1}}
\begin{itemize}
\item Antes: pasaporte + directivas mínimas; supervisión y rescates nacionales. 2008: rescates descoordinados; Fortis partido. 2009 \autor{de Larosière} $\to$ 2011 \concepto{SESF} (ABE, AEVM, AESPJ, JERS) \vertema{3.B.36}.
\item 2010-2012: crisis de deuda; España, 555 pb, hasta 100.000 millones (41.000 del MEDE) \vertema{4.A.18}; fragmentación.
\item 26/6 y 5/12/2012: \concepto{Informe de los cuatro presidentes} (\autor{Van Rompuy}). 29/6/2012: Cumbre del Euro: «romper el círculo vicioso». 26/7/2012: «lo que sea necesario» y OMT.
\item Reglamento del MUS (2013) $\to$ BCE supervisa desde el 4/11/2014; MUR desde 2016; \concepto{Informe de los cinco presidentes} (22/6/2015).
\end{itemize}
\graficorepaso[\linewidth]{cronologia}{Unión bancaria (arriba) y mercados de capitales (abajo), 2009-2027}

\sub{I.2. ¿Por qué una unión bancaria?\label{rep:I.2}}
\textbf{Círculo vicioso en forma reducida} \vertema{3.B.26}
\begin{itemize}
\item \emph{Supuestos}: Estado: $s=s_0+\alpha L$; bancos con deuda propia: $L=L_0+\beta s$; $\alpha\beta<1$.
\item \emph{Desarrollo}:
\begin{formula}
\funcionclave{s^*=\frac{s_0+\alpha L_0}{1-\alpha\beta}}
\end{formula}
$\Delta s^*=\frac{\alpha}{1-\alpha\beta}\Delta L_0$: multiplicador $\frac{1}{1-\alpha\beta}>1$.
\item \emph{Implicaciones}: resolución, FUR y EDIS bajan $\alpha$; supervisión baja $L_0$; deuda pública con capital o activo seguro bajan $\beta$. \emph{Extensiones}: abrazo mortal (\autor{Farhi}-\autor{Tirole}, 2018); ESBies. \emph{Evidencia}: deuda pública 9,3\,\% del activo; parte doméstica 39\,\% (2014) $\to$ 28\,\% (2025).
\end{itemize}
\graficorepaso[0.85\linewidth]{bucle-banco-soberano}{Perturbación $\Delta L_0$: de $E$ a $E'$, más que el impulso; con unión bancaria ($\alpha\downarrow$), $E''$}

\textbf{Trilema financiero} (\autor{Schoenmaker}, 2011)
\begin{itemize}
\item Estabilidad, integración y políticas nacionales: solo 2. Banco transfronterizo: salvarlo cuesta $C$, evita $B$; el país de origen recibe $\theta B$:
\begin{formula}
\funcionclave{\theta B\geq C\ \text{(nacional)}\quad\text{vs}\quad B\geq C}
\end{formula}
Si $\theta<C/B<1$: infraprovisión; crece con la integración. Salidas: \emph{ring-fencing}, inestabilidad (euro hasta 2012) o unión bancaria. Captura y tolerancia \vertema{3.A.20}. Quien decide debe pagar.
\item \emph{Evidencia}: Fortis (2008) partido por países; liquidez retenida en las filiales (2012). \emph{Valoración}: argumento sencillo y potente, pero exige compartir también los costes.
\end{itemize}
\graficorepaso[0.75\linewidth]{trilema-financiero}{Solo 2 de 3: la unión bancaria renuncia a las políticas nacionales}

\sub{I.3. El código normativo único\label{rep:I.3}}
\begin{itemize}
\item Consejo Europeo, 2009: reglamentos, pocas opciones nacionales; los 27.
\item \textbf{CRR/CRD} (Basilea~III \vertema{3.B.27}); \textbf{CRR3/CRD~VI}: desde el 1/1/2025, \emph{output floor} 50\,\% $\to$ 72,5\,\% (2030), transitorios a 2032; CRD~VI transpuesta en enero de 2026; terceros países 2027; FRTB desde 2027.
\item \textbf{BRRD} (2014/59): planes, resolución, \emph{bail-in}. \textbf{DGSD} (2014/49): 100.000\,€, 7 días, 0,8\,\% (2024). BRRD + MUR + DGSD = \concepto{CMDI}.
\item \textbf{ABE} (2011, París): normas técnicas, directrices, código interactivo, pruebas de resistencia.
\end{itemize}

% ==========================================================================
\bloque[10,5]{II. Los pilares de la unión bancaria\label{rep:II}}
\begin{ideaclave}
Supervisión única: funciona. Resolución única: reglas y fondo lleno, poco uso y respaldo bloqueado. Garantía común: nunca aprobada, retirada en 2026. Una casa con 2 pilares y medio.
\end{ideaclave}
\graficorepaso[0.95\linewidth]{pilares-union-bancaria}{Los 3 pilares sobre el código normativo único; el tercero, sin construir}

\sub{II.1. Mecanismo Único de Supervisión\label{rep:II.1}}
\begin{itemize}
\item Art.~127.6 TFUE; \textbf{Reglamento 1024/2013}; BCE + ANC; 21 países (Bulgaria en el euro desde 2026; cooperación estrecha 2020 de Bulgaria y Croacia).
\item 109 \concepto{significativas} ($>$30.000 millones, $>$20\,\% PIB, 3 mayores, transfronteriza, ayuda del MEDE), $\approx$85\,\% de los activos, con equipos conjuntos; 1.828 menos significativas (77\,\% en DE y AT). Licencias, para todas.
\item Consejo de Supervisión (\autor{Claudia Buch}, 2024) $\to$ Consejo de Gobierno por no objeción; separación de la política monetaria.
\item \concepto{PRES}: modelo de negocio, gobierno, capital, liquidez $\to$ P2R y P2G. Para 2026: CET1 11,2\,\%, total 15,6\,\%, P2R 1,2\,\% (105 bancos). Reforma plurianual.
\item Evaluación global 2014: 130 bancos, déficit de 24.600 millones en 25.
\end{itemize}
\graficorepaso[\linewidth]{solvencia-morosidad}{CET1 del 12,7\,\% (2015) al 16,0\,\% (2026); dudosos del 7,5\,\% al 1,9\,\%}

\sub{II.2. Mecanismo Único de Resolución\label{rep:II.2}}
\begin{itemize}
\item \textbf{Reglamento 806/2014}; JUR (Bruselas): 115 bancos, 92 con resolución prevista. Inviable (BCE) $\to$ sin alternativa privada $\to$ interés público $\to$ dispositivo (JUR, Comisión, Consejo en 24 h). Si no, liquidación nacional.
\item Instrumentos: venta, banco puente, segregación, \concepto{bail-in}.
\item \concepto{FUR}: 1\,\% de los depósitos garantizados; $>$81.000 millones (2025); sin aportaciones 2024-2026; acuerdo intergubernamental de 2014, mutualizado en 2024. Solo tras el 8\,\% de \emph{bail-in}; máximo 5\,\%.
\item \textbf{Respaldo del MEDE}: 68.000 millones; firmado el 27/1/2021; Italia no ratifica (21/12/2023) \vertema{3.B.45}.
\end{itemize}
\graficorepaso[0.9\linewidth]{fondo-unico-resolucion}{FUR: de 10.800 millones (2016) a más de 81.000 (2025)}

\textbf{\emph{Bail-in} frente a \emph{bail-out}}
\begin{itemize}
\item \emph{Supuestos}: quiebra con probabilidad $\pi$; recuperación $R$; acreedores neutrales al riesgo. Con rescate, $r=r_f$. Con \emph{bail-in}:
\begin{formula}
\funcionclave{r-r_f\approx\pi(1-R)}
\end{formula}
= subvención implícita del demasiado grande para caer.
\item \emph{Implicaciones}: inconsistencia temporal $\Rightarrow$ MREL (objetivo medio 27,8\,\% de APR; déficit 200 millones) y TLAC. \emph{Extensiones}: contagio; Credit Suisse (AT1, 2023). \emph{Evidencia}: subvención hasta 300.000 millones de dólares en el euro (FMI, 2012); CDS tras Chipre (\autor{Schäfer}, \autor{Schnabel} y \autor{Weder di Mauro}, 2016).
\end{itemize}
\graficorepaso[0.9\linewidth]{jerarquia-acreedores}{Primero el capital, después la deuda; los depósitos, con preferencia; los garantizados, nunca}

\textbf{Casos}: Popular (7/6/2017, Santander, 1\,€); Vicenza y Veneto (2017, liquidación nacional con ayudas); ABLV (2018), PNB (2019); Sberbank (2022). \textbf{CMDI}: acuerdo 25/6/2025; Directivas 2026/806 y 2026/804 y Reglamento 2026/808 (abril de 2026), aplicación en mayo de 2028: más interés público, FGD financian transferencias (cuentan para el 8\,\%), preferencia de los depositantes en 3 niveles.

\sub{II.3. La garantía común de depósitos\label{rep:II.3}}
\begin{itemize}
\item \concepto{EDIS} (24/11/2015): reaseguro, coaseguro, seguro pleno en 2024. Bloqueo: reducir riesgos antes o compartirlos. 2017: 2 fases; Eurogrupo 2022: solo CMDI. Julio de 2026: se sustituirá; programa de trabajo 2027 (4/11/2026): retirada; nuevas propuestas en el 1.\textsuperscript{er} trimestre de 2027.
\end{itemize}
\textbf{Seguro común y riesgo moral}
\begin{itemize}
\item Garantía creíble elimina el pánico de \autor{Diamond}-\autor{Dybvig} \vertema{3.B.26}. $n$ sistemas, desviación $\sigma$, correlación $\rho$:
\begin{formula}
\funcionclave{\frac{F_c}{F_n}=\sqrt{\rho+\frac{1-\rho}{n}}<1}
\end{formula}
Riesgo moral: la garantía es una \emph{put} (\autor{Merton}, 1977), y entre países. Remedio: aportaciones según riesgo, supervisión común, reducir riesgos heredados.
\item \emph{Evidencia}: con 0,8\,\%, suficiente incluso en crisis grave y sin subvenciones cruzadas sistemáticas (\autor{Carmassi} \emph{et al.}, 2018).
\end{itemize}
\graficorepaso[0.85\linewidth]{seguro-depositos-comun}{Fondo común frente a nacional: menor con más países y menos correlación}
\begin{itemize}
\item \textbf{Deuda pública}: pondera 0 y sin límite de grandes riesgos; Basilea (2017): sin consenso; reforma gradual + activo seguro \vertema{3.B.37}.
\end{itemize}

% ==========================================================================
\bloque[8,5]{III. La Unión de Mercados de Capitales\label{rep:III}}
\begin{ideaclave}
Mucha financiación bancaria y ahorro en depósitos: los riesgos se reparten mal. Diez años de normas (UMC) y pocos cambios de fondo; en 2025, Unión de Ahorros e Inversiones.
\end{ideaclave}

\sub{III.1. ¿Por qué?\label{rep:III.1}}
\begin{itemize}
\item Préstamos 73\,\% de la deuda de las empresas (EE.\,UU.: 37\,\%; bancos 12\,\%); capitalización 71\,\% del PIB (EE.\,UU. 224\,\%); depósitos 31\,\% de los activos de los hogares (EE.\,UU. 14\,\%); 300.000 millones al año salen (\autor{Letta}).
\item Rueda de repuesto, innovación, inversión (\autor{Draghi}: 4,4-4,7\,\% del PIB), reparto de riesgos.
\end{itemize}
\graficorepaso[\linewidth]{capitalizacion-bursatil}{Capitalización bursátil: UE 71\,\%, EE.\,UU. 224\,\% y Japón 172\,\% del PIB (2025)}
\graficorepaso[\linewidth]{financiacion-empresas}{Préstamos sobre deuda de las empresas: 73\,\% en el euro, 37\,\% en EE.\,UU.}
\textbf{Reparto de riesgos} (\autor{Asdrubali}, \autor{Sørensen} y \autor{Yosha}, 1996; \autor{Ingram} \vertema{3.B.16})
\begin{itemize}
\item PIB $\to$ renta nacional (capital) $\to$ renta disponible (fiscal) $\to$ consumo (crédito):
\begin{formula}
\funcionclave{1=\beta_K+\beta_F+\beta_C+\beta_U}
\end{formula}
EE.\,UU.: 39/13/23, sin suavizar 25\,\%. Euro (1997-2022): capital $\approx$10, crédito $\approx$30, fiscal $\approx$0, sin suavizar $\approx$60\,\%. El capital es el canal más estable; el crédito se esfuma en las crisis.
\end{itemize}
\graficorepaso[\linewidth]{reparto-riesgos}{Lo que suaviza cada canal: en el euro, poco el capital y casi nada el presupuesto}

\sub{III.2. La UMC (2015-2024)\label{rep:III.2}}
\begin{itemize}
\item Plan de 30/9/2015 (\autor{Juncker}, más de 30 acciones); revisión 2017 (Brexit); plan de 24/9/2020: 16 acciones, 3 objetivos.
\item Logros: STS (2017/2402, 2019); folleto y \emph{Listing Act} (2024/2809); ELTIF (2023); \emph{crowdfunding}; ESAP (julio de 2027); PEPP (2022; 2 proveedores); inversión minorista (acuerdo 12/2025); cinta consolidada (acciones, septiembre de 2026); EMIR~3 (cuenta activa); T+1 (11/10/2027); FASTER (2030); insolvencia (2026).
\item Balance: normas sí; integración, poca.
\end{itemize}

\sub{III.3. La Unión de Ahorros e Inversiones\label{rep:III.3}}
\begin{itemize}
\item \autor{Letta} (4/2024) y \autor{Draghi} (9/2024). Estrategia de 19/3/2025: ciudadanos y ahorro; inversión y financiación; integración y escala; supervisión.
\item Cuentas de ahorro e inversión (30/9/2025); titulización (acuerdo 21/10/2026); pensiones complementarias (11/2025) \vertema{4.B.24}; integración y supervisión (4/12/2025): AEVM supervisa plataformas significativas, ECC, DCV y criptoactivos (Consejo, 9/10/2026); 28.º régimen (3/2026).
\item Euro digital: propuesta 2023; Consejo 12/2025; PE 7/2026; trílogos; emisión posible 2029 \vertema{3.B.43}.
\item Valoración: razón política nueva y supervisión europea sobre la mesa; lo difícil depende de los Estados.
\end{itemize}

% ==========================================================================
\bloque[1]{Conclusión}
\begin{itemize}
\item \textbf{Recapitulación}: 2012, círculo vicioso y trilema; código común; MUS y MUR con FUR; falta EDIS. UMC y UAI: el otro motor.
\item \textbf{Relevancia}: transmisión monetaria, crisis bancarias sin crisis de deuda, financiar la transición.
\item \textbf{Extensiones}: \vertema{3.B.27} \vertema{3.B.26} \vertema{3.B.16} \vertema{3.B.43} \vertema{3.B.45} \vertema{3.B.39} \vertema{4.A.19} \vertema{4.A.20}.
\item \textbf{Opinión}: garantía común con menos sesgo doméstico; AEVM para los mercados grandes.
\item \textbf{Idea final}: el euro empezó la casa por el tejado; desde 2012 se levantan los cimientos. La casa estará terminada cuando el seguro de los depósitos sostenga el tejado.
\end{itemize}

\bloque{Datos clave}
\begin{itemize}
\item 109 significativas (1/9/2026), $\approx$85\,\% de activos; 21 países; CET1 16,0\,\%, dudosos 1,89\,\% (2,17\,\% sin efectivo), ROE 10,7\,\% (2T 2026).
\item FUR $>$81.000 millones; respaldo 68.000; MREL 27,8\,\%.
\item Capitalización: UE 71\,\%, EE.\,UU. 224\,\% del PIB (2025).
\end{itemize}

\bloque{Conceptos}
\begin{itemize}
\item \textbf{Inviable o probablemente inviable}: incumple la autorización, activos menores que pasivos, no puede pagar o necesita ayuda pública extraordinaria.
\item \textbf{Interés público}: la resolución mantiene funciones esenciales, evita el contagio y protege fondos públicos y depositantes mejor que la liquidación.
\item \textbf{Ningún acreedor peor parado}: nadie pierde en la resolución más que en una liquidación.
\item \textbf{Recapitalización precautoria}: ayuda a un banco solvente por el déficit de una prueba de resistencia (Monte dei Paschi, 2017).
\item \textbf{Cinta consolidada}: un único flujo de precios y volúmenes de todas las plataformas.
\item \textbf{Requisitos}: Pilar~1 CET1 4,5\,\%, nivel~1 6\,\%, total 8\,\%; apalancamiento 3\,\%; conservación 2,5\,\%; TLAC 18\,\% de APR y 6,75\,\% de la exposición.
\end{itemize}

\bloque{Autores y fechas}
\begin{itemize}
\item 1983 \autor{Diamond}-\autor{Dybvig}; 1977 \autor{Merton}; 1996 \autor{Asdrubali}-\autor{Sørensen}-\autor{Yosha}; 2009 \autor{de Larosière}; 2011 \autor{Schoenmaker}; 2012 4 presidentes; 2014 MUS; 2015 5 presidentes, EDIS, UMC; 2016 MUR; 2016 \autor{Brunnermeier} \emph{et al.}; 2017 Popular; 2018 \autor{Farhi}-\autor{Tirole}; 2020 UMC 2.0; 2021 MEDE; 2024 \autor{Letta}, \autor{Draghi}; 2025 CRR3, UAI; 2026 CMDI, EDIS retirado.
\end{itemize}

\bloque{Preguntas del tribunal}
\begin{itemize}
\item \textbf{¿Por qué la supervisión sola no basta?} Quien decide debe pagar: sin resolución, fondo y garantía comunes, $\alpha$ sigue alto.
\item \textbf{¿Es creíble el \emph{bail-in}?} Solo con MREL suficiente en manos de quien pueda perder y un respaldo común.
\item \textbf{¿Por qué no se aprueba el EDIS?} Riesgos heredados y deuda pública en los balances: reducir antes o compartir antes.
\item \textbf{¿Por qué Vicenza y Veneto no se resolvieron?} Sin interés público según la JUR: liquidación italiana con ayudas; la CMDI amplía esa prueba.
\item \textbf{¿Qué aporta la AEVM como supervisor?} Lo mismo que el MUS a los bancos: un solo supervisor para los mercados transfronterizos grandes, contra el trilema y la competencia entre supervisores.
\item \textbf{¿Basta el reparto de riesgos privado?} Reduce la necesidad de presupuesto común, pero el crédito se retira en las crisis: hace falta también un respaldo público.
\item \textbf{¿Y España?} BdE en el MUS (sus grandes bancos, en el BCE), FROB en el MUR, FGD; CET1 del 13,5\,\%, el más bajo; Popular (2017); MARF y BME Growth \vertema{4.A.19} \vertema{4.A.20}.
\item \textbf{¿Mercados o bancos?} Complementarios: los mercados reparten riesgos entre países y financian la innovación; los bancos, a las pymes.
\end{itemize}

\bloque{Trampas del test}
\begin{itemize}
\item El código normativo único es de los 27; la unión bancaria, del euro y la cooperación estrecha.
\item El FUR actúa tras el 8\,\% de \emph{bail-in}; máximo 5\,\%. La garantía: 100.000\,€ por depositante y \emph{entidad}.
\item El MUS supervisa \emph{directamente} solo las significativas; licencias, todas.
\item El respaldo del MEDE no está en vigor; el EDIS nunca se aprobó.
\item El MUS empezó en noviembre de 2014; el MUR, plenamente, en enero de 2016.
\item Cooperación estrecha: Bulgaria y Croacia, desde el 1/10/2020; ahora, ninguna.
\item La CMDI no crea EDIS: reforma BRRD, MUR y DGSD.
\end{itemize}

\end{columnas}
\end{document}
